\chapter[The Type A2 Quantum Enveloping Algebra]{The Type \(A_2\) Quantum Enveloping Algebra}
\label{ch:quantum-enveloping}

We construct the two-root quantum enveloping algebra and its formal
classical limit. The transition from one simple root to two requires
the quantum Serre relations as well as the Cartan and mixed
commutator relations. The rank-one algebra, its Hopf structure,
and its finite-dimensional type-1 representations are treated
in Results H10--H15 of \emph{Topics in Rings, Modules, and
Associative Algebras}~\cite[PDF pp.~25--28]{mg-rm03d}.
Here the main tasks are to prove a basis for the two-root algebra,
check the coproduct on the new relations, and compute its
first-order skew part. The constructions below use
\(\Delta^{\mathrm{op}}\) for tensor factor exchange applied to
\(\Delta\), and \(a\wedge b=a\otimes b-b\otimes a\).

\section{The quantum Serre relations and the positive part}

Let \(A\) be a commutative coefficient ring and \(q\in A^\times\).
Write \(c=q+q^{-1}\) and put
\[
 \mathcal S(x,y)=x^2y-cxyx+yx^2 .
\]

\begin{theorem}[The two-generator positive part]
\label{res:UQ01}
The \(A\)-algebra on \(x,y\) with relations
\(\mathcal S(x,y)=\mathcal S(y,x)=0\) is free as an
\(A\)-module with basis \(x^a z^b y^d\), where
\[
 z=xy-q^{-1}yx,\qquad a,b,d\in\mathbb Z_{\ge0}.
\]
In this algebra \(xz=qzx\), \(zy=qyz\), and \(yx=qxy-qz\).
This basis and presentation commute with arbitrary base change
on \(A\) carrying \(q\) to a unit.
\end{theorem}
\begin{proof}
The expressions \(xz-qzx\) and \(zy-qyz\) expand to
\(\mathcal S(x,y)\) and \(\mathcal S(y,x)\), respectively.
The remaining relation is the defining formula for \(z\)
solved for \(yx\). Conversely these three relations for
\(x,z,y\) eliminate \(z\) and give precisely the two Serre
relations.

Adjoin \(z\) to \(A[x]\) with \(zx=q^{-1}xz\).
On the resulting algebra let \(\sigma(x)=qx\) and
\(\sigma(z)=q^{-1}z\). This preserves the relation and
has inverse given by reciprocal scalars. Define the
\(\sigma\)-derivation by \(\delta(x)=-qz\), \(\delta(z)=0\),
where \(\delta(ab)=\sigma(a)\delta(b)+\delta(a)b\).
It respects the relation because both \(\delta(zx)\) and
\(\delta(q^{-1}xz)\) are \(-z^2\). Hence adjoining \(y\)
with \(ya=\sigma(a)y+\delta(a)\) gives the second pair
of commutation rules. The Ore normal form in
\emph{Topics in Rings, Modules, and Associative Algebras},
Result O11~\cite[PDF p.~34]{mg-rm03d}, makes the indicated
monomials a free basis. Every relation and each basis
applies over the new coefficient ring by the same construction.
\end{proof}

At \(q=1\), the generator \(z\) is \([x,y]\) and commutes with
\(x,y\). In the matrix realization of \(\mathfrak{sl}_3\), the
assignment \(x\mapsto e_{12}\), \(y\mapsto e_{23}\),
\(z\mapsto e_{13}\) then identifies this presentation
with the enveloping algebra of the positive nilpotent subalgebra.
The commutation rules and the basis on each side prove the
identification, using the ordered PBW basis of
\emph{Lie Algebras and Finite-Dimensional Representations},
Result A05~\cite[PDF pp.~8--9]{mg-lie}, on the classical side.

\section{The Cartan coefficients and triangular normal form}

The common Cartan matrix will be \(C=(a_{ij})_{1\le i,j\le 2}\),
with \(a_{11}=a_{22}=2\) and \(a_{12}=a_{21}=-1\).
We give parallel coefficient choices so that the normal-form
argument can also be used to justify the formal limit.

\begin{definition}[The generic and truncated formal presentations]
\label{res:UQ02}
Fix a field \(k\) of characteristic zero. The generic choice is
\(A=k\) with a parameter \(q\in k^\times\), \(q^2\ne1\), or
\(A=k(q)\) with \(q\) indeterminate. It uses
\[
 T=A[K_1^{\pm1},K_2^{\pm1}],\qquad
 \tau_i(K_l)=q^{a_{li}}K_l,\qquad
 D_i=\frac{K_i-K_i^{-1}}{q-q^{-1}} .
\]
The other choice, for an integer \(N\ge1\), uses
\[
 \begin{gathered}
 A=k[h]/(h^N),\qquad T=A[H_1,H_2],\qquad
 q=e^h,\qquad K_i=e^{hH_i},\\
 \tau_i(H_l)=H_l+a_{li},\qquad
 D_i=\frac{\sinh(hH_i)/h}{\sinh(h)/h}\pmod {h^N}.
 \end{gathered}
\]
The divided series are first formed in \(k[H_1,H_2][[h]]\);
the denominator there has constant term \(1\).
Both choices have \(\tau_i(K_l)=q^{a_{li}}K_l\), and the
\(\tau_i\) commute. The basic relations on \(T\) and \(E_i,F_i\) are
\[
 t F_j=F_j\tau_j^{-1}(t),\qquad
 E_i t=\tau_i^{-1}(t)E_i,\qquad
 [E_i,F_j]=\delta_{ij}D_i \quad (t\in T).
\]
The quantum enveloping algebra is obtained by also imposing
\(\mathcal S(E_i,E_j)=\mathcal S(F_i,F_j)=0\) for \(i\ne j\).
Write \(U\) for the generic choice and \(U_N\) for the truncated
formal choice.
\end{definition}

\begin{lemma}[Normal form before the Serre quotient]
\label{res:UQ03}
With the coefficients of Definition~\ref{res:UQ02},
let \(B_0\) impose only the basic relations on \(T,E_i,F_i\).
Multiplication identifies \(B_0\), as an \(A\)-module, with
\[
 A\langle F_1,F_2\rangle\otimes_A T
       \otimes_A A\langle E_1,E_2\rangle .
\]
There is an algebra involution \(\omega\) interchanging
\(E_i,F_i\) and inverting the Cartan generators
(\(K_i\mapsto K_i^{-1}\) or \(H_i\mapsto-H_i\)).
There is also an algebra anti-involution \(\eta\) interchanging
\(E_i,F_i\) and fixing \(T\).
\end{lemma}
\begin{proof}
Use the basic relations to move all negative letters left of
the \(T\)-coefficient and all positive letters right of it,
and combine adjacent \(T\)-coefficients. To see termination,
first count the total number of \(E,F\) letters, then the
total number of factors, and then inversions relative to the
order negative, Cartan, positive. Replacing \(E_iF_j\) by its
commutator term lowers the first count. Combining adjacent
Cartan factors lowers the second. Every remaining interchange
lowers the inversion count. One can start with Cartan factors
from a monomial basis of \(T\) and expand all replacements
linearly.

Disjoint replacements commute. Overlaps with adjacent
Cartan factors are compatible with multiplication since the
\(\tau_i\) are algebra automorphisms. The only remaining
overlap is \(E_i tF_j\). In the two orders the leading term
is \(F_j\tau_i^{-1}\tau_j^{-1}(t)E_i\); the other term in
one order is \(\delta_{ij}\tau_i^{-1}(t)D_i\) and in
the other is \(\delta_{ij}D_i\tau_j^{-1}(t)\).
These coincide by commutativity of the translations and
of \(T\). The normal form is therefore independent of reductions:
induct on the ordered counts and compare the first two choices,
using the overlap computation or commuting them when disjoint.
Normalizing all terms in a defining relation inserted into
a word gives zero, so the normal form descends to the
presented algebra and proves independence as well as spanning.

Cartan inversion \(\kappa\) satisfies
\(\kappa\tau_i=\tau_i^{-1}\kappa\) and \(\kappa(D_i)=-D_i\).
These identities check the two Cartan and mixed relations
after the swap defining \(\omega\). Reversing all products
while interchanging \(E,F\) gives the Cartan relations in
the opposite order, and sends
\([E_i,F_j]\) to \([E_j,F_i]\). This checks \(\eta\).
\end{proof}

\begin{lemma}[Crossing the Serre ideal]
\label{res:UQ04}
In the algebra \(B_0\) of Lemma~\ref{res:UQ03},
\(\mathcal S(E_i,E_j)\) commutes with every \(F_l\), and
\(\mathcal S(F_i,F_j)\) commutes with every \(E_l\).
Under the same normal form, the two-sided ideal generated
by the positive and negative Serre polynomials is the sum
of the submodules obtained by putting the corresponding
free-algebra Serre ideal in a positive or negative factor.
\end{lemma}
\begin{proof}
Put \(S=\mathcal S(E_i,E_j)\), \(i\ne j\).
For \(l=j\) use
\(
 [S,F_j]=(\tau_i^{-2}D_j-c\tau_i^{-1}D_j+D_j)E_i^2.
\)
For \(l=i\), the commutator rule
\([ab,F]=a[b,F]+[a,F]b\) gives
\[
\begin{split}
 [S,F_i]
 &=\bigl(D_i+\tau_i^{-1}D_i-c\tau_i^{-1}\tau_j^{-1}D_i\bigr)E_iE_j\\
 &\quad+\bigl(\tau_j^{-1}(D_i+\tau_i^{-1}D_i)-cD_i\bigr)E_jE_i .
\end{split}
\]
Every displayed coefficient is zero. In the generic case,
clear \(q-q^{-1}\) and evaluate each coefficient on
the terms \(K_l^s\) with \(s=1,-1\).
The three scalar identities are
\[
 q^{2s}-c q^s+1=0,\qquad
 1+q^{-2s}-c q^{-s}=0,\qquad
 q^s(1+q^{-2s})-c=0 .
\]
They follow directly from \(c=q+q^{-1}\).
In the formal case do the same calculation on
\(K_l^s=e^{shH_l}\) before reduction.
The relations with \(D_l\) then follow by cancellation
of \(e^h-e^{-h}\) in the formal coefficient domain,
so they hold in every truncation. Applying \(\omega\)
gives the negative commutators.

To check the ideal statement, take the submodule with
positive Serre ideal in the positive factor. Left
multiplication by \(F_l\) and right multiplication by \(E_l\)
preserve it at once. Cartan factors shift across a word
according to the total number of its root letters; a Serre
polynomial is homogeneous for those counts.
Thus left or right Cartan multiplication preserves
the submodule. Moving a left \(E_l\) across a negative
word uses only the basic relations and produces either
a remaining positive letter or shorter negative words
with Cartan coefficients. This also preserves the positive
ideal factor. For right multiplication by \(F_l\), it
suffices to move it across a positive expression \(pSq\).
The commutators \([p,F_l]\) and \([q,F_l]\) have only
Cartan coefficients and positive words. The commutator
\([S,F_l]\) is zero, and moving a Cartan coefficient
across \(S\) preserves its Serre factor by homogeneity.
This proves closure on both sides. Applying \(\eta\) gives
the same result for the negative submodule without changing
the factor order. Their sum is a two-sided ideal containing
the generators, and all its basis-type terms already
belong to the generated ideal.
\end{proof}

\begin{theorem}[Triangular PBW in type \(A_2\)]
\label{res:UQ05}
For either coefficient choice of Definition~\ref{res:UQ02},
multiplication is an \(A\)-module isomorphism from
the negative Serre algebra tensored with \(T\) and then
with the positive Serre algebra onto the quantum
enveloping algebra. In particular, writing
\(z_-=F_1F_2-q^{-1}F_2F_1\) and
\(z_+=E_1E_2-q^{-1}E_2E_1\), a basis consists of
\[
 F_1^a z_-^b F_2^d\, T_0\, E_1^{a'} z_+^{b'} E_2^{d'},
\]
where the six root exponents are nonnegative integers and
\(T_0\) ranges over the indicated Cartan monomials.
In the generic case the subalgebra generated by
\(F_i,K_i^{\pm1},E_i\) is the rank-one algebra
\(U_q(\mathfrak{sl}_2)\).
\end{theorem}
\begin{proof}
Apply Lemmas~\ref{res:UQ03} and \ref{res:UQ04};
passing to the Serre quotient in either factor gives
the asserted tensor product by right exactness of
tensoring. Theorem~\ref{res:UQ01} supplies the bases
of the two Serre factors, and \(T\) has its evident
monomial basis. The fixed-index generators satisfy the
rank-one relations, and the corresponding monomials
\(F_i^aK_i^bE_i^d\) are a subset of the independent
displayed family. This proves that the rank-one map is
injective by the rank-one PBW result H10
~\cite[PDF p.~25]{mg-rm03d}.
\end{proof}

\section{The Hopf structure}

\begin{theorem}[Coproduct and antipode in type \(A_2\)]
\label{res:UQ06}
The formulas
\[
\begin{aligned}
 \Delta(E_i)&=E_i\otimes1+K_i\otimes E_i, &
 \Delta(F_i)&=F_i\otimes K_i^{-1}+1\otimes F_i,\\
 \epsilon(E_i)&=\epsilon(F_i)=0, &
 S(E_i)&=-K_i^{-1}E_i,\qquad S(F_i)=-F_iK_i
\end{aligned}
\]
give Hopf structures on \(U\) and \(U_N\) from
Definition~\ref{res:UQ02}. In the generic case \(K_i\)
is group-like with counit \(1\) and antipode \(K_i^{-1}\).
In the formal case \(H_i\) is primitive with counit
zero and antipode \(-H_i\), so \(K_i=e^{hH_i}\) is
group-like. The antipodes are bijective, and the formal
structures commute with reduction in \(N\).
\end{theorem}
\begin{proof}
The proposed Cartan coproduct preserves its own relations.
For the generic case, conjugation by \(K_l\otimes K_l\)
scales the images of \(E_i,F_i\) by \(q^{a_{li}},q^{-a_{li}}\).
For the formal case, commuting with \(H_l\otimes1+1\otimes H_l\)
gives the same two weights in additive form. Thus the
Cartan-root relations hold. In the commutator of
\(\Delta(E_i),\Delta(F_j)\) the mixed term
\([K_i\otimes E_i,F_j\otimes K_j^{-1}]\) vanishes:
the factors needed to interchange the two products
are \(q^{-a_{ij}}\) and \(q^{a_{ji}}\), which cancel.
The term \([E_i\otimes1,1\otimes F_j]\) also vanishes.
The other terms are therefore
\[
 \delta_{ij}\bigl(D_i\otimes K_i^{-1}+K_i\otimes D_i\bigr)
       =\delta_{ij}\Delta(D_i).
\]
For Laurent polynomials this last identity follows from
multiplying by \(q-q^{-1}\) and using \(\Delta(K_i)=K_i\otimes K_i\).
In the formal case it follows by the same calculation
before cancellation of \(e^h-e^{-h}\) in the
formal series ring in the two Cartan sets, then reduction.

For the positive Serre relation take \(i\ne j\) and put
\[
 x=E_i\otimes1,\quad u=K_i\otimes E_i,\qquad
 y=E_j\otimes1,\quad v=K_j\otimes E_j.
\]
We need \(\mathcal S(x+u,y+v)=0\). The pure parts are
\(\mathcal S(x,y)\) and \(\mathcal S(u,v)\); the latter
is zero because the first-factor \(K\)'s commute and
the second factor is the Serre polynomial.
Use \(ux=q^2xu,\ uy=q^{-1}yu,\ vx=q^{-1}xv\).
Terms with \(y\) and two copies of \(u\), or with \(v\) and two
copies of \(x\), have scalar factor \(1-cq^{-1}+q^{-2}=0\).
Terms with \(y\) and one copy each of \(x,u\) sum to
\[
 cxyu-c(xyu+qyxu)+(1+q^2)yxu=0.
\]
Those with \(v\) and one copy each of \(x,u\) sum to
\[
 (1+q^2)xuv-c(xvu+qxuv)+cxvu=0 .
\]
This covers every term. The negative relation follows
by \(\omega\) from Lemma~\ref{res:UQ03} on both factors
and then tensor-factor exchange. Those transformations
send the displayed positive coproduct formula to the
negative one. Hence \(\Delta\) is an algebra map.

The counit respects the relations, including \(\epsilon(D_i)=0\).
Coassociativity follows on the generators: for \(E_i\) each
iterate is \(E_i\otimes1\otimes1+K_i\otimes E_i\otimes1+
K_i\otimes K_i\otimes E_i\); the \(F_i\) expansion similarly
puts \(K_i^{-1}\) after each \(F_i\). The formulas for group-like
and primitive elements are coassociative directly. The counit
identities also hold on the generators.

To check that \(S\) extends anti-multiplicatively, use the
Cartan inversion \(\kappa\) of Lemma~\ref{res:UQ03}.
The identity \(\kappa\tau_i=\tau_i^{-1}\kappa\) checks
the reversed Cartan-root relations, since \(S(E_i)\)
and \(S(F_i)\) have their respective original weights.
For the mixed relation, move \(K_i^{-1}\) to the left
and \(K_j\) to the right. Symmetry of \(a_{ij}\) gives
\[
 [S(F_j),S(E_i)]
    =K_i^{-1}(F_jE_i-E_iF_j)K_j=-\delta_{ij}D_i,
\]
which is the required value of \(S(\delta_{ij}D_i)\).
For each reversed word in a positive Serre polynomial,
move all \(K^{-1}\)'s to the left. The scalar in addition
to the three minus signs is the product of \(q^{a_{rs}}\)
over the three pairs of letters. This is independent of
their order because \(a_{rs}=a_{sr}\). Here the pair sum
is \(2-1-1=0\). The result is the common Cartan monomial
times the reversed Serre polynomial, which is zero.
The negative check moves all \(K\)'s to the right,
with the reciprocal scalar, and gives the same conclusion.

The two antipode products give the counit times \(1\)
on every generator by substitution. They extend to products
by anti-multiplicativity: for instance in
\(\sum S((ab)_{(1)})(ab)_{(2)}\) the two contributions of
\(a\) occur next to one another, giving \(\epsilon(a)\),
and then the contributions of \(b\) give \(\epsilon(b)\).
For the other product the contributions of \(b\) occur
first. Finally \(S^2\) fixes \(T\) and scales \(E_i,F_i\)
by \(q^{-2},q^2\), respectively. These reciprocal scalings
preserve the defining relations and have inverses.
Thus \(S^2\) is invertible and \(S\) is bijective.
All formulas in the formal choice use reductions of the
same series, so they commute with reduction in \(N\).
\end{proof}

\section{The formal classical limit}

\begin{theorem}[A topologically free quantization of the enveloping algebra]
\label{res:UQ07}
Use the compatible truncations of Definition~\ref{res:UQ02}
and set \(U_h=\varprojlim_N U_N\).
It is a topologically free \(k[[h]]\)-module carrying the
completed Hopf structure of Theorem~\ref{res:UQ06}.
Its reduction modulo \(h\) is \(U(\mathfrak{sl}_3(k))\),
as a Hopf algebra. Here the completed tensor product
is the inverse limit of \(U_N\otimes_{k[h]/(h^N)} U_N\).
\end{theorem}
\begin{proof}
The monomials in Theorem~\ref{res:UQ05} have compatible
reductions. Thus each \(U_N\) is the free module on the
same index set over \(k[h]/(h^N)\), and reduction is
coefficientwise. An element of the limit is a series
\(\sum_{n\ge0}h^n v_n\), where each \(v_n\) is a finite
\(k\)-linear combination of the indexed monomials.
Indeed, a finite truncation has finite support and
conversely a finite list of such \(v_n\)'s has finite
total support. This proves topological freeness and
identifies the quotient by \(h\) with \(U_1\).
The same statement with pairs of basis indices proves
topological freeness of the completed tensor product.
The Hopf maps are the compatible maps in the theorem
cited above, so their identities pass to inverse limits.

Modulo \(h\), the Cartan commutators in Definition~\ref{res:UQ02}
give \([H_l,E_i]=a_{li}E_i\), \([H_l,F_i]=-a_{li}F_i\),
and \([E_i,F_j]=\delta_{ij}H_i\). The Serre relations
become the ordinary double commutator relations.
The assignment to \(\mathfrak{sl}_3(k)\) sends
\[
 E_1\mapsto e_{12},\quad E_2\mapsto e_{23},\quad
 F_1\mapsto e_{21},\quad F_2\mapsto e_{32},\quad
 H_i\mapsto e_{ii}-e_{i+1,i+1},
\]
where \(e_{ij}\) is a matrix unit. These matrices satisfy
the displayed relations by matrix multiplication.
The resulting algebra map to \(U(\mathfrak{sl}_3(k))\)
sends \(z_+,z_-\) to \(e_{13},-e_{31}\).
Its basis maps precisely to the ordered PBW monomials
for this basis of \(\mathfrak{sl}_3(k)\), by Result A05 of
\emph{Lie Algebras and Finite-Dimensional Representations}
~\cite[PDF pp.~8--9]{mg-lie}. Thus the map is an
isomorphism. The coproduct reduces to the primitive formula
on the indicated generators, and the counit and antipode
reduce to the standard ones. This also proves the
Hopf assertion.
\end{proof}

\begin{theorem}[The first-order cobracket on \(\mathfrak{sl}_3\)]
\label{res:UQ08}
On the specialization of Theorem~\ref{res:UQ07}, use
lowercase \(e_i,f_i,h_i\) for the images of \(E_i,F_i,H_i\).
For a lift \(X\in U_h\) of \(x\), put
\[
 \delta(x)=\frac{\Delta(X)-\Delta^{\mathrm{op}}(X)}h
             \pmod h .
\]
Its restriction to \(\mathfrak{sl}_3(k)\) is a Lie
bialgebra cobracket. We identify \(x\wedge y\) with
\(x\otimes y-y\otimes x\). Explicitly,
\[
\begin{aligned}
 \delta(h_i)&=0, &
 \delta(e_i)&=h_i\wedge e_i, &
 \delta(f_i)&=h_i\wedge f_i,\\
 \delta([e_1,e_2])&=(h_1+h_2)\wedge[e_1,e_2]+2e_1\wedge e_2,\\
 \delta([f_1,f_2])&=(h_1+h_2)\wedge[f_1,f_2]-2f_1\wedge f_2 .
\end{aligned}
\]
\end{theorem}
\begin{proof}
The specialization is cocommutative, so the numerator
is divisible by \(h\). The completed tensor product
is \(h\)-torsion-free by Theorem~\ref{res:UQ07}.
Replacing a lift by \(X+hY\) does not change the
answer, since the specialization of
\(\Delta(Y)-\Delta^{\mathrm{op}}(Y)\) is zero.
Expanding \(K_i=1+hH_i+O(h^2)\) and
\(K_i^{-1}=1-hH_i+O(h^2)\) in Theorem~\ref{res:UQ06}
gives the displayed simple-root and Cartan values.

Write \(\Delta_0\) for the specialized coproduct.
Multiplicativity of \(\Delta,\Delta^{\mathrm{op}}\)
and division of their skew difference give
\[
 \delta([x,y])=[\Delta_0(x),\delta(y)]
                  +[\delta(x),\Delta_0(y)] .
\]
For primitive \(x,y\) the right-hand side is
\(\operatorname{ad}_x\delta(y)-\operatorname{ad}_y\delta(x)\)
for the adjoint action on both factors. Using the
two off-diagonal Cartan entries \(-1\) gives
the two remaining displayed formulas. The eight
matrix basis elements from the preceding proof are thus
sent into the skew part of \(\mathfrak{sl}_3\otimes\mathfrak{sl}_3\).
The commutator identity gives the Lie-algebra cocycle rule
on that subspace.

Here is the co-Jacobi check. Write \(B=\Delta-\Delta^{\mathrm{op}}\)
and \(D=(\Delta\otimes1)\Delta=(1\otimes\Delta)\Delta\).
In tensor notation, write \(D\) as \(123\) and denote
its factor permutations by their orders. Expansion gives
\((B\otimes1)B=123-213-231+321\).
The sum of its three cyclic permutes is zero, because
\(123,231\) are cyclic permutations and \(213,321\) are
also cyclic permutations. The same coefficientwise
description of the triple completed tensor product
shows it is \(h\)-torsion-free. Dividing by \(h^2\)
and reducing gives the co-Jacobi identity for \(\delta\).
Pairing this cyclic identity with triples of linear forms
gives the Jacobi identity for the dual bracket. Together
with the skew and cocycle rules this proves the Lie
bialgebra assertion.
\end{proof}

\begin{corollary}[The coboundary tensor in the matrix model]
\label{res:UQ09}
With the convention of Theorem~\ref{res:UQ08},
the same cobracket on \(\mathfrak{sl}_3\) is
\(\delta(x)=\operatorname{ad}_x(r)\), where
\[
 r=-e_{12}\wedge e_{21}-e_{23}\wedge e_{32}-e_{13}\wedge e_{31}.
\]
\end{corollary}
\begin{proof}
The matrix-unit rule is
\([e_{ij},e_{lm}]=\delta_{jl}e_{im}-\delta_{mi}e_{lj}\).
Each wedge in \(r\) has total weight zero, giving
\(\operatorname{ad}_{h_i}(r)=0\).
For \(e_1=e_{12}\), the first wedge with its minus sign
contributes \(h_1\wedge e_1\); the other two contributions
are \(-e_{13}\wedge e_{32}\) and \(e_{13}\wedge e_{32}\).
The same computation for \(e_2=e_{23}\) gives
\(h_2\wedge e_2\), with cross contributions
\(e_{13}\wedge e_{21}\) and \(-e_{13}\wedge e_{21}\).
The Lie algebra automorphism \(x\mapsto-x^{\mathsf T}\)
sends \(r\) to \(-r\), \(e_i\) to \(-f_i\) and
\(h_i\) to \(-h_i\). Applying it to the identities for
\(e_i\) gives \(\operatorname{ad}_{f_i}(r)=h_i\wedge f_i\).
Both cobrackets satisfy the cocycle identity, for
\(\operatorname{ad}_x(r)\) by the Jacobi identity.
Their equality on the six generators therefore
implies equality on all of \(\mathfrak{sl}_3\).
\end{proof}

\paragraph{Sources.}
Kassel~\cite[Sections XVI.5 and XVII.2]{kassel} gives the
general Drinfeld--Jimbo constructions. Etingof--Semenyakin
~\cite[Sections 4.1--4.2]{etingof-semenyakin} treats the
Serre presentation and quantization in general.
This chapter proves the full normal form and formal limit
for type \(A_2\) directly, with a specified skew-coproduct
convention.
